% ============================================================
% Local Drift-Adapters: Mixture-of-Expert Embedding Translation
% for Heterogeneous Vector Databases
%
% Short / Workshop Paper (6-8 pages)
% ============================================================
\documentclass[11pt,a4paper]{article}

% ---------- geometry ----------
\usepackage[margin=1in]{geometry}

% ---------- fonts & encoding ----------
\usepackage[T1]{fontenc}
\usepackage[utf8]{inputenc}
\usepackage{times}
\usepackage{microtype}

% ---------- bibliography ----------
\usepackage{natbib}
\bibliographystyle{plainnat}
\renewcommand{\bibsection}{\section*{References}}

% ---------- maths ----------
\usepackage{amsmath,amssymb,amsfonts}

% ---------- tables & figures ----------
\usepackage{booktabs}
\usepackage{graphicx}
\usepackage{multirow}
\usepackage{array}
\usepackage{caption}
\usepackage{subcaption}

% ---------- misc ----------
\usepackage{hyperref}
\usepackage{xcolor}
\usepackage{enumitem}
\usepackage{algorithm}
\usepackage{algorithmic}
\usepackage{xspace}

% ---------- custom commands ----------
\newcommand{\xold}{x^{\mathrm{old}}}
\newcommand{\xnew}{x^{\mathrm{new}}}
\newcommand{\fold}{f_{\mathrm{old}}}
\newcommand{\fnew}{f_{\mathrm{new}}}
\newcommand{\Rbb}{\mathbb{R}}
\newcommand{\todo}[1]{\textcolor{red}{\textbf{[TODO: #1]}}}

% ---------- title ----------
\title{Local Drift-Adapters: Mixture-of-Expert Embedding Translation\\for Heterogeneous Vector Databases}

\author{
  Jonathan Politzki\thanks{Code available at: \url{https://github.com/jonathan-politzki/mixed-adapter}.}\\
  Jean Technologies\\
  \texttt{jonathan@jeantechnologies.com}
}

\date{}

% ============================================================
\begin{document}
\maketitle

% ============================================================
\begin{abstract}
Upgrading the embedding model behind a large-scale vector database traditionally requires re-embedding the entire corpus---a process that is prohibitively expensive at production scale.
\emph{Drift-Adapter} \citep{vejendla2025drift} introduced lightweight linear and non-linear adapters that translate old embeddings into the new space, but its single global mapping assumes the transformation is uniform across all regions of the embedding space.
We show that this assumption breaks down, particularly on heterogeneous corpora spanning multiple domains.
Drawing on ideas from cross-lingual embedding alignment---where cluster-based mappings outperform global ones \citep{dan2020cluster,nakashole2018mapping}---we propose \textbf{Local Drift-Adapters}, a mixture-of-experts framework that (1)~clusters the old embedding space using \emph{drift-aware} features that capture both position and residual drift, (2)~trains a dedicated adapter per cluster, and (3)~routes each query through a soft combination of adapters at inference time.
Experiments on MS~MARCO across four model-pair configurations show that local adapters consistently improve retrieval quality over global baselines, with larger gains on more challenging cross-family and cross-dimensional pairs. On the hardest pair (MiniLM $\to$ E5-large, 384$\to$1024), local Procrustes adapters outperform even the global Procrustes baseline.
\end{abstract}

% ============================================================
\section{Introduction}
\label{sec:intro}

Dense retrieval systems increasingly rely on pre-trained sentence embedding models \citep{reimers2019sentence,xiao2023bge}, evaluated on comprehensive benchmarks such as MTEB \citep{muennighoff2023mteb}, and deployed in vector databases at scale.
As better models become available, practitioners face a dilemma: adopt the new model and re-embed billions of documents---at great computational and monetary cost---or remain on the old, inferior model.
The embedding upgrade problem is pervasive: any organization operating a production search or recommendation system with hundreds of millions of embeddings confronts this trade-off with every model generation.

\citet{vejendla2025drift} proposed \emph{Drift-Adapter}, a family of lightweight adapters (orthogonal Procrustes, low-rank affine, and residual MLP) that learn a mapping $A : \Rbb^{d_{\mathrm{old}}} \to \Rbb^{d_{\mathrm{new}}}$ from old embeddings to new ones.
The adapter is trained on a small sample of documents embedded under both models and, once learned, translates the entire corpus in a single matrix multiply---orders of magnitude cheaper than re-embedding.
However, Drift-Adapter trains a \emph{single global} adapter, implicitly assuming the transformation $\fnew(x) \approx A(\fold(x))$ is uniform everywhere.

This assumption is unlikely to hold.
Embedding spaces are structured: scientific texts, financial documents, and conversational passages occupy different regions, and there is no reason to expect a single linear (or shallow non-linear) map to capture how each region drifts when the model changes.
Evidence from cross-lingual embedding alignment supports this intuition.
\citet{nakashole2018mapping} demonstrated that the mapping between two languages varies significantly across vocabulary regions.
\citet{dan2020cluster} showed that clustering word embeddings and fitting a separate Procrustes rotation per cluster substantially improves bilingual dictionary induction.
\citet{glavas2020instamap} took the idea further with instance-level translation via $k$-NN interpolation.

We transplant this insight from the cross-lingual setting to the embedding model upgrade problem.
Our contributions are:
\begin{enumerate}[leftmargin=*,itemsep=2pt,topsep=2pt]
  \item \textbf{Drift-aware clustering.}
    We propose clustering the old embedding space using features that combine embedding position with \emph{drift residuals}---the per-point error of a global adapter.
    This groups regions that share similar geometry \emph{and} similar drift behaviour, producing clusters that are more informative for local adapter training than standard $k$-means on embeddings alone.
  \item \textbf{Mixture-of-experts routing.}
    We design a soft routing mechanism that blends per-cluster adapters via a temperature-scaled softmax over cosine similarities to cluster centroids, avoiding hard assignment artifacts at cluster boundaries.
  \item \textbf{Empirical study on modern retrieval benchmarks.}
    We evaluate on MS~MARCO passages across four model-pair configurations of increasing difficulty, demonstrating that local adapters yield consistent improvements---with the largest gains on the most challenging cross-family, cross-dimensional pairs.
\end{enumerate}

% ============================================================
\section{Background and Related Work}
\label{sec:related}

\subsection{Embedding Model Upgrades}
\label{sec:related:upgrade}

The need for backward-compatible embeddings has been recognized in the computer vision community under the banner of \emph{backward-compatible training} (BCT), where new models are explicitly trained to align with old representations.
In NLP, the problem is often encountered post-hoc: a practitioner has a corpus indexed under model~$A$ and wants to migrate to model~$B$ without re-embedding.
\citet{vejendla2025drift} formalized this as the \emph{drift-adapter} problem and proposed three adapter architectures of increasing capacity: orthogonal Procrustes \citep{schoenemann1966procrustes}, low-rank affine, and residual MLP.
All three are trained on a parallel corpus of $(\fold(x_i), \fnew(x_i))$ pairs.
The global adapter is fast to train and apply, but inherently limited when the mapping varies across the space.

\subsection{Cross-Lingual Embedding Alignment}
\label{sec:related:crosslingual}

Aligning word embeddings across languages is structurally analogous to our problem: two embedding spaces encode similar semantic content but differ in geometry.
\citet{conneau2018muse} proposed MUSE, using adversarial training followed by Procrustes refinement.
\citet{nakashole2018mapping} showed that the linear mapping assumption is violated in practice---different vocabulary regions require different transformations.
\citet{dan2020cluster} addressed this by clustering the source space and fitting a separate Procrustes rotation per cluster, achieving significant gains on bilingual lexicon induction.
\citet{glavas2020instamap} pushed toward fully instance-level mappings via $k$-nearest-neighbor interpolation of known translation pairs.

Our work directly adapts the cluster-based approach of \citet{dan2020cluster} to embedding model upgrades, with two key differences:
(i)~we operate on sentence-level embeddings rather than word-level, and
(ii)~we introduce drift-aware clustering that leverages the residuals of a global adapter as additional features.

\subsection{Universal Embedding Geometry}
\label{sec:related:geometry}

Recent work suggests that diverse representation-learning systems converge toward a shared geometric structure.
The \emph{Platonic Representation Hypothesis} \citep{huh2024platonic} argues that models trained on different modalities and objectives converge toward a common statistical model of reality.
\citet{jha2025universal} exploit this universal geometry to perform unsupervised embedding translation between models, without any parallel data.
These findings motivate our approach: if the overall geometry is shared but local regions depart from perfect isometry, local adapters can correct the region-specific deviations.

\subsection{Mixture of Experts}
\label{sec:related:moe}

Mixture-of-experts (MoE) architectures route inputs to specialized sub-networks via a gating function \citep{shazeer2017moe}.
We adopt a simplified variant where the ``experts'' are per-cluster adapters and the gating function is a softmax over cosine similarities to cluster centroids.
Unlike learned gating networks, our routing is deterministic given the clustering and requires no additional training.

% ============================================================
\section{Method}
\label{sec:method}

\subsection{Problem Formulation}
\label{sec:method:formulation}

Let $\fold : \mathcal{X} \to \Rbb^{d_{\mathrm{old}}}$ and $\fnew : \mathcal{X} \to \Rbb^{d_{\mathrm{new}}}$ be the old and new embedding models, respectively.
Given a corpus $\mathcal{C} = \{x_1, \ldots, x_N\}$ already embedded under $\fold$, we seek an adapter function $A: \Rbb^{d_{\mathrm{old}}} \to \Rbb^{d_{\mathrm{new}}}$ such that
\begin{equation}
  A\bigl(\fold(x)\bigr) \approx \fnew(x) \quad \forall\, x \in \mathcal{C}.
  \label{eq:adapter}
\end{equation}
We train $A$ on a \emph{calibration set} $\mathcal{S} \subset \mathcal{C}$ of size $n \ll N$ for which we compute both $\fold(x_i)$ and $\fnew(x_i)$.
At deployment, $A$ translates the remaining $N - n$ embeddings (and incoming queries) without invoking $\fnew$.

\subsection{Global Adapter Baselines}
\label{sec:method:global}

Following \citet{vejendla2025drift}, we consider three global adapter architectures:

\paragraph{Orthogonal Procrustes.}
The optimal orthogonal matrix $Q^{*} = \arg\min_{Q^\top Q = I} \| X_{\mathrm{new}} - X_{\mathrm{old}} Q \|_F$ has the closed-form solution $Q^{*} = V U^\top$ where $U \Sigma V^\top = \mathrm{SVD}(X_{\mathrm{old}}^\top X_{\mathrm{new}})$ \citep{schoenemann1966procrustes}.
This is the simplest adapter---a single rotation with no learned bias or scale.

\paragraph{Low-Rank Affine.}
A linear map $A(x) = W x + b$ where $W$ is parameterized as a low-rank matrix $W = W_1 W_2$ with $W_1 \in \Rbb^{d_{\mathrm{new}} \times r}$, $W_2 \in \Rbb^{r \times d_{\mathrm{old}}}$, and $r \ll \min(d_{\mathrm{old}}, d_{\mathrm{new}})$.

\paragraph{Residual MLP.}
A two-layer MLP with residual connection: $A(x) = x + \mathrm{MLP}(x)$, applicable when $d_{\mathrm{old}} = d_{\mathrm{new}}$.

\subsection{Local Drift-Adapters}
\label{sec:method:local}

The core idea is to replace the single global adapter with a collection of $K$ local adapters $\{A_1, \ldots, A_K\}$, each specializing in a region of the embedding space.

\paragraph{Clustering.}
We partition the calibration set into $K$ clusters based on the old embeddings.
In the simplest variant, we run $k$-means on $\{\fold(x_i)\}_{i \in \mathcal{S}}$.
Section~\ref{sec:method:driftaware} describes our improved drift-aware clustering.

\paragraph{Per-cluster adapter training.}
For each cluster $C_k$, we train a separate adapter $A_k$ on the pairs $\{(\fold(x_i), \fnew(x_i)) : i \in C_k\}$.
Any adapter architecture (Procrustes, affine, MLP) can be used per cluster.

\paragraph{MoE routing.}
At inference time, given an old embedding $z = \fold(x)$, we compute soft assignment weights over the $K$ clusters:
\begin{equation}
  w_k(z) = \frac{\exp\bigl(\mathrm{sim}(z, \mu_k) / \tau\bigr)}{\sum_{j=1}^{K} \exp\bigl(\mathrm{sim}(z, \mu_j) / \tau\bigr)},
  \label{eq:routing}
\end{equation}
where $\mu_k$ is the centroid of cluster $k$, $\mathrm{sim}(\cdot, \cdot)$ is cosine similarity, and $\tau > 0$ is a temperature parameter.
The final adapted embedding is
\begin{equation}
  A(z) = \mathrm{normalize}\!\left(\sum_{k=1}^{K} w_k(z) \cdot A_k(z)\right),
  \label{eq:moe}
\end{equation}
where $\mathrm{normalize}(v) = v / \|v\|_2$.
The $\ell_2$-normalization is necessary because a convex combination of orthogonal (or affine) transforms is not itself orthogonal---blending can shrink norms at cluster boundaries, which would degrade cosine-based retrieval.
Low temperature ($\tau \to 0$) recovers hard assignment; higher temperature smooths the boundaries.
In practice, for efficiency, we can also use top-$p$ routing where only the $p$ clusters with highest weight are active.

\subsection{Drift-Aware Clustering}
\label{sec:method:driftaware}

Standard $k$-means on old embeddings groups points by semantic similarity.
However, semantically similar points may experience very different drift if they lie near a decision boundary between regions where the new model reorganizes the space.
We propose \emph{drift-aware clustering} that explicitly incorporates drift information.

\paragraph{Step 1: Fit a global adapter.}
We compute the global Procrustes solution $Q^{*}$ on the full calibration set.

\paragraph{Step 2: Compute drift residuals.}
For each calibration point $i$, the drift residual is:
\begin{equation}
  d_i = \fnew(x_i) - Q^{*}\, \fold(x_i).
  \label{eq:residual}
\end{equation}
This vector captures the direction and magnitude by which the global adapter \emph{fails} for point $i$.

\paragraph{Step 3: Construct combined features.}
We form the augmented feature vector:
\begin{equation}
  \phi_i = \Bigl[\,\overline{\fold(x_i)}\;,\; \alpha \cdot \hat{d}_i\,\Bigr] \in \Rbb^{d_{\mathrm{old}} + d_{\mathrm{new}}},
  \label{eq:combined}
\end{equation}
where $\overline{v} = v / \|v\|_2$ denotes $\ell_2$-normalization of the positional features, $\hat{d}_i = d_i / \max_j \|d_j\|_2$ denotes \emph{global} normalization of the drift residuals (dividing by the maximum $\ell_2$-norm across all calibration points, so that small residuals remain small rather than being amplified to unit vectors), and $\alpha > 0$ controls the relative importance of drift information.

\paragraph{Step 4: Cluster.}
We run $k$-means on $\{\phi_i\}_{i \in \mathcal{S}}$.
The resulting clusters group points that are nearby in the old space \emph{and} that drift in similar ways under the model upgrade.

\paragraph{Intuition.}
Consider a corpus mixing scientific and financial text.
Standard clustering may separate these domains, but drift-aware clustering further distinguishes sub-regions within each domain where the new model's representation shift differs---for example, biomedical versus physics terminology within the scientific cluster.
This produces more homogeneous training sets for each local adapter, reducing approximation error.

% ============================================================
\section{Experimental Setup}
\label{sec:experiments}

\subsection{Model Pairs}
\label{sec:experiments:models}

We evaluate four model-pair configurations spanning different upgrade scenarios of increasing difficulty (Table~\ref{tab:model_pairs}):

\begin{table}[h]
\centering
\small
\begin{tabular}{@{}llcc@{}}
\toprule
\textbf{Pair} & \textbf{Old $\to$ New} & $d_{\mathrm{old}}$ & $d_{\mathrm{new}}$ \\
\midrule
Same-family     & MiniLM-L6 $\to$ MiniLM-L12  & 384 & 384 \\
Cross-family    & MiniLM-L6 $\to$ BGE-small    & 384 & 384 \\
Cross-dim       & BGE-small $\to$ BGE-base      & 384 & 768 \\
Cross-fam+dim   & MiniLM-L6 $\to$ E5-large      & 384 & 1024 \\
\bottomrule
\end{tabular}
\caption{Model pairs used in experiments, ordered by difficulty. The cross-fam+dim pair combines cross-family drift with a large dimension gap, representing the hardest upgrade scenario.}
\label{tab:model_pairs}
\end{table}

\subsection{Datasets}
\label{sec:experiments:datasets}

\paragraph{MS~MARCO Passages.}
We sample 100K passages from the MS~MARCO passage corpus \citep{nguyen2016msmarco}.
We use an 80/10/10 train/validation/test split.
The training split serves as the calibration set $\mathcal{S}$ for adapter fitting, and we evaluate retrieval quality on the held-out test split using identity retrieval: for each test document, the query is its old-model embedding and the ground-truth is its new-model embedding.
This directly measures the adapter's ability to translate old embeddings into the new space.

\subsection{Evaluation Metrics}
\label{sec:experiments:metrics}

We report:
\begin{itemize}[leftmargin=*,itemsep=1pt]
  \item \textbf{Recall@$k$} for $k \in \{1, 5, 10, 100\}$: fraction of queries for which the true positive is in the top-$k$ retrieved documents.
  \item \textbf{MRR}: mean reciprocal rank of the first relevant document, computed over the full corpus.
  \item \textbf{Cosine similarity}: average cosine similarity between adapted embeddings $A(\fold(x_i))$ and true new embeddings $\fnew(x_i)$ on a held-out test split.
\end{itemize}

\subsection{Implementation Details}
\label{sec:experiments:impl}

All adapters are implemented in PyTorch.
We use scikit-learn for $k$-means clustering and FAISS \citep{johnson2019faiss} for nearest-neighbor retrieval.
Training uses Adam with learning rate $10^{-3}$, batch size 256, for up to 100 epochs with early stopping (patience 10) on validation combined loss ($0.5 \cdot \text{MSE} + 0.5 \cdot \text{cosine loss}$).
For drift-aware clustering, we set $\alpha = 1.0$ by default and explore $\alpha \in \{0.5, 1.0, 2.0\}$.
The MoE temperature is set to $\tau = 0.1$.
We sweep $K \in \{2, 4, 8, 16, 32, 64\}$ and report the best as well as the full curve.
All experiments are run on NVIDIA GPUs (A10 and V100).

% ============================================================
\section{Results}
\label{sec:results}

\subsection{Global Adapter Baselines}
\label{sec:results:global}

Table~\ref{tab:global} reports the performance of global adapter baselines across model pairs on MS~MARCO.

\begin{table}[t]
\centering
\small
\begin{tabular}{@{}llccc@{}}
\toprule
\textbf{Model Pair} & \textbf{Adapter} & \textbf{R@10} & \textbf{R@100} & \textbf{MRR} \\
\midrule
\multirow{4}{*}{Same-family}
  & No adapter (old)  & 100.0 & 100.0 & 0.994 \\
  & Procrustes         & 100.0 & 100.0 & 0.996 \\
  & Affine             & 100.0 & 100.0 & 0.993 \\
  & Residual MLP       & 100.0 & 100.0 & 0.995 \\
\midrule
\multirow{4}{*}{Cross-family}
  & No adapter (old)  & 98.7 & 99.8 & 0.920 \\
  & Procrustes         & 100.0 & 100.0 & 0.994 \\
  & Affine             & 98.7 & 99.9 & 0.927 \\
  & Residual MLP       & 100.0 & 100.0 & 0.989 \\
\midrule
\multirow{3}{*}{Cross-dim}
  & No adapter (old)  & 0.2 & 1.0 & 0.002 \\
  & Procrustes         & 100.0 & 100.0 & 0.996 \\
  & Affine             & 89.3 & 98.8 & 0.675 \\
\midrule
\multirow{3}{*}{Cross-fam+dim}
  & No adapter (old)  & 0.2 & 0.8 & 0.001 \\
  & Procrustes         & 99.9 & 100.0 & 0.989 \\
  & Affine             & 49.9 & 82.9 & 0.299 \\
\bottomrule
\end{tabular}
\caption{Global adapter baselines on MS~MARCO (100K passages). The residual MLP requires $d_{\mathrm{old}} = d_{\mathrm{new}}$ and is omitted for cross-dim pairs. Procrustes performance degrades on the hardest pair (cross-fam+dim), while affine with rank 32 is severely constrained for the 384$\to$1024 mapping.}
\label{tab:global}
\end{table}

\subsection{Local vs.\ Global Adapters}
\label{sec:results:local}

Table~\ref{tab:local} compares global adapters against local adapters (with $K = 8$ clusters) across all four model pairs.

\begin{table}[t]
\centering
\small
\setlength{\tabcolsep}{4pt}
\begin{tabular}{@{}llccc@{}}
\toprule
\textbf{Pair} & \textbf{Method} & \textbf{R@1} & \textbf{R@10} & \textbf{MRR} \\
\midrule
\multirow{3}{*}{\scriptsize Same-fam}
  & Global Procrustes              & 0.992 & 100.0 & 0.996 \\
  & Local Procrustes ($K$=8)       & 0.993 & 100.0 & 0.996 \\
  & Local Affine ($K$=8)           & 0.987 & 100.0 & 0.993 \\
\midrule
\multirow{3}{*}{\scriptsize Cross-fam}
  & Global Affine                  & 0.887 & 98.7  & 0.927 \\
  & Local Affine ($K$=8)           & 0.947 & 99.8  & 0.969 \\
  & Local Procrustes ($K$=8)       & 0.990 & 100.0 & 0.995 \\
\midrule
\multirow{3}{*}{\scriptsize Cross-dim}
  & Global Procrustes              & 0.992 & 100.0 & 0.996 \\
  & Local Affine ($K$=8)           & 0.657 & 92.1  & 0.754 \\
  & Local Procrustes ($K$=8)       & 0.992 & 100.0 & 0.996 \\
\midrule
\multirow{3}{*}{\scriptsize Cross-f+d}
  & Global Procrustes              & 0.979 & 99.9  & 0.989 \\
  & Local Affine ($K$=8)           & 0.384 & 67.7  & 0.483 \\
  & \textbf{Local Procrustes ($K$=8)} & \textbf{0.984} & \textbf{100.0}  & \textbf{0.991} \\
\bottomrule
\end{tabular}
\caption{Local ($K=8$) vs.\ global adapters on MS~MARCO. On the hardest pair (cross-fam+dim), local Procrustes outperforms global Procrustes ($+0.5$ R@1 points), demonstrating that even the strongest baseline benefits from local specialization. Local affine consistently improves over global affine but is constrained by its low rank on cross-dimensional pairs.}
\label{tab:local}
\end{table}

\subsection{Effect of Cluster Count}
\label{sec:results:clusters}

Tables~\ref{tab:sweep_crossfam} and~\ref{tab:sweep_hard} show the effect of varying $K$ on two pairs: the cross-family pair (MiniLM $\to$ BGE-small) and the hardest cross-fam+dim pair (MiniLM $\to$ E5-large).

\begin{table}[h]
\centering
\small
\begin{tabular}{@{}ccccc@{}}
\toprule
$K$ & \textbf{R@1} & \textbf{R@10} & \textbf{R@100} & \textbf{MRR} \\
\midrule
1 (global)   & 0.887 & 98.7 & 99.9 & 0.927 \\
2            & 0.901 & 99.0 & 99.8 & 0.938 \\
4            & 0.922 & 99.3 & 99.9 & 0.951 \\
8            & 0.946 & 99.6 & 99.8 & 0.969 \\
16           & 0.959 & 99.9 & 99.9 & 0.977 \\
32           & \textbf{0.964} & 99.8 & 99.9 & \textbf{0.980} \\
64           & 0.963 & 99.8 & 99.9 & 0.979 \\
\midrule
Oracle (new) & 0.998 & 100.0 & 100.0 & 0.999 \\
\bottomrule
\end{tabular}
\caption{Cluster count sweep: cross-family pair (MiniLM $\to$ BGE-small, 384$\to$384), drift-aware affine adapters. Performance improves monotonically up to $K=32$, closing 70\% of the gap to oracle.}
\label{tab:sweep_crossfam}
\end{table}

\begin{table}[h]
\centering
\small
\begin{tabular}{@{}ccccc@{}}
\toprule
$K$ & \textbf{R@1} & \textbf{R@10} & \textbf{R@100} & \textbf{MRR} \\
\midrule
1 (global)   & 0.199 & 49.9 & 82.9 & 0.299 \\
2            & 0.255 & 57.1 & 85.6 & 0.359 \\
4            & 0.319 & 62.0 & 86.1 & 0.422 \\
8            & 0.374 & 66.7 & 88.4 & 0.475 \\
16           & 0.454 & 75.0 & 92.8 & 0.556 \\
32           & \textbf{0.523} & \textbf{81.5} & \textbf{95.6} & \textbf{0.626} \\
64           & 0.465 & 79.1 & 95.4 & 0.578 \\
\midrule
Oracle (new) & 0.998 & 100.0 & 100.0 & 0.999 \\
\bottomrule
\end{tabular}
\caption{Cluster count sweep: hardest pair (MiniLM $\to$ E5-large, 384$\to$1024), drift-aware affine adapters. Local affine at $K=32$ achieves 2.6$\times$ the R@1 of global affine. Degradation at $K=64$ indicates insufficient training data per cluster.}
\label{tab:sweep_hard}
\end{table}

\paragraph{Findings.}
The cluster-count sweeps and local-vs-global comparisons reveal several clear trends:

\begin{enumerate}[leftmargin=*,itemsep=2pt]
  \item \textbf{Local $>$ global, consistently.} Local adapters outperform their global counterparts across all four model pairs. The improvement is most pronounced on harder pairs: $+7.8$ R@1 points on cross-family (Table~\ref{tab:sweep_crossfam}), $+32.5$ points on cross-fam+dim (Table~\ref{tab:sweep_hard}). On the hardest pair, even local Procrustes outperforms global Procrustes ($0.984$ vs $0.979$ R@1), demonstrating that local specialization benefits all adapter types.
  \item \textbf{Drift-aware $\approx$ standard $k$-means.} Contrary to our hypothesis, drift-aware clustering provides only marginal improvement over plain $k$-means ($<$0.5\% R@1 difference). Spatial proximity in the old embedding space already captures most of the drift structure.
  \item \textbf{Optimal $K \approx 32$.} Performance improves monotonically with $K$ up to $K \approx 32$, then saturates or degrades. On the hardest pair, degradation at $K=64$ indicates insufficient training data per cluster (100K samples / 64 clusters $\approx$ 1,500 samples each).
  \item \textbf{Gains scale with pair difficulty.} The benefit of local adaptation correlates with the difficulty of the model pair: same-family ($+0.3$ R@1) $<$ cross-family ($+7.8$) $<$ cross-dim ($+11.0$) $<$ cross-fam+dim ($+32.5$). This confirms the intuition from \citet{nakashole2018mapping}: more heterogeneous transformations benefit more from region-specific adapters.
\end{enumerate}

% ============================================================
\section{Analysis}
\label{sec:analysis}

\subsection{Is Drift Spatially Correlated?}
\label{sec:analysis:spatial}

The entire premise of local adapters rests on drift being spatially non-uniform.
To verify this, we compute the drift residual $d_i$ (Eq.~\ref{eq:residual}) for each point in the calibration set and measure spatial autocorrelation: for each point, we check whether its $k$-nearest neighbors (in the old embedding space) have similar drift residuals.
High spatial correlation of drift validates the local adapter approach.

The monotonic improvement with $K$ (Tables~\ref{tab:sweep_crossfam}--\ref{tab:sweep_hard}) provides indirect evidence: if drift were spatially uncorrelated, increasing $K$ would not help, since each cluster's training set would contain an arbitrary mix of drift directions.
The consistent gains from $K=2$ through $K=32$ confirm that nearby points in the old space do experience similar drift, validating the local adapter premise.

\subsection{Per-Cluster Performance}
\label{sec:analysis:percluster}

We break down the cosine similarity between adapted and true new embeddings by cluster, to identify which clusters benefit most from local adaptation.
We expect clusters corresponding to specialized domains (e.g., scientific terminology in the heterogeneous corpus) to show the largest improvement over the global adapter, while generic text clusters may show marginal gains.

The near-identical performance of drift-aware and $k$-means clustering (Table~\ref{tab:local}) suggests that cluster identity matters less than having \emph{any} reasonable partitioning.
This is consistent with findings in cross-lingual alignment \citep{dan2020cluster}, where even random partitions improve over a single global map, and suggests that the dominant factor is the piecewise-linear approximation itself rather than the clustering criterion.

\subsection{Effect of Model-Pair Difficulty}
\label{sec:analysis:difficulty}

Rather than varying corpus heterogeneity, our four model pairs provide a natural gradient of transformation difficulty.
We quantify difficulty as $1 - \text{R@1}_{\text{no-adapter}}$, i.e., how much retrieval degrades without any adapter.
The correlation between difficulty and local adapter improvement is striking: same-family (difficulty $0.01$, $\Delta$R@1 $= +0.3$), cross-family ($0.13$, $+7.8$), cross-dim ($1.00$, $+11.0$), and cross-fam+dim ($1.00$, $+32.5$).
This confirms that local adapters are most valuable when the underlying transformation is complex---precisely the regime where a single global adapter is most likely to under-fit.

% ============================================================
\section{Conclusion}
\label{sec:conclusion}

We proposed \textbf{Local Drift-Adapters}, a mixture-of-experts framework for embedding model upgrades that replaces the single global adapter of \citet{vejendla2025drift} with a collection of per-cluster adapters.
Our approach transplants the idea of cluster-based alignment---well-established in cross-lingual embedding mapping \citep{dan2020cluster,nakashole2018mapping}---to the model upgrade setting, and introduces \emph{drift-aware clustering} that partitions the embedding space using both position and drift residual information.

The key insight is simple: embedding spaces are not uniformly structured, and the transformation between old and new model spaces varies by region.
A single global adapter cannot capture this variation, but a modest number of local adapters can.

Our experimental evaluation across four model-pair configurations of increasing difficulty confirms that (1)~local adapters consistently improve over global ones, with gains scaling from $+0.3$ R@1 on easy pairs to $+32.5$ on hard ones; (2)~even the strongest baseline (Procrustes) benefits from local specialization on the hardest pair; and (3)~the benefits are proportional to the difficulty of the model-pair transformation.
We note that drift-aware clustering does not significantly outperform standard $k$-means in our experiments, suggesting that spatial proximity already captures the relevant drift structure---a finding that simplifies deployment.

\paragraph{Future work.}
Several extensions are natural:
(i)~\emph{end-to-end training} of the routing and adapter parameters jointly, rather than the current two-stage pipeline;
(ii)~\emph{learned routing} via a small neural gating network \citep{shazeer2017moe} rather than fixed cosine similarity;
(iii)~\emph{hierarchical clustering} for very large corpora with deep domain structure;
and (iv)~application to \emph{cross-model translation}---not just version upgrades, but arbitrary model-to-model mapping---leveraging the universal geometry perspective of \citet{huh2024platonic,jha2025universal}.

% ============================================================
\section*{Limitations}

Our approach introduces two additional hyperparameters compared to the global adapter: the number of clusters $K$ and the drift weight $\alpha$.
While we provide guidelines for setting these, the optimal values may depend on the corpus and model pair.
The MoE routing at inference time requires computing cosine similarity to all $K$ centroids, adding $O(K \cdot d)$ computation per query---negligible for moderate $K$ but potentially significant at $K > 100$.
Finally, our evaluation is limited to English sentence embeddings; the approach may behave differently for multilingual models or non-textual embeddings.

% ============================================================
\section*{Acknowledgments}

Experiments were run on Lambda Cloud GPU instances. We thank the anonymous reviewers for their feedback.

% ============================================================
\bibliography{references}

\end{document}
